% Options for packages loaded elsewhere
\PassOptionsToPackage{unicode}{hyperref}
\PassOptionsToPackage{hyphens}{url}
\documentclass[
  ignorenonframetext,
]{beamer}
\newif\ifbibliography
\usepackage{pgfpages}
\setbeamertemplate{caption}[numbered]
\setbeamertemplate{caption label separator}{: }
\setbeamercolor{caption name}{fg=normal text.fg}
\beamertemplatenavigationsymbolsempty
% remove section numbering
\setbeamertemplate{part page}{
  \centering
  \begin{beamercolorbox}[sep=16pt,center]{part title}
    \usebeamerfont{part title}\insertpart\par
  \end{beamercolorbox}
}
\setbeamertemplate{section page}{
  \centering
  \begin{beamercolorbox}[sep=12pt,center]{section title}
    \usebeamerfont{section title}\insertsection\par
  \end{beamercolorbox}
}
\setbeamertemplate{subsection page}{
  \centering
  \begin{beamercolorbox}[sep=8pt,center]{subsection title}
    \usebeamerfont{subsection title}\insertsubsection\par
  \end{beamercolorbox}
}
% Prevent slide breaks in the middle of a paragraph
\widowpenalties 1 10000
\raggedbottom
\AtBeginPart{
  \frame{\partpage}
}
\AtBeginSection{
  \ifbibliography
  \else
    \frame{\sectionpage}
  \fi
}
\AtBeginSubsection{
  \frame{\subsectionpage}
}
\usepackage{iftex}
\ifPDFTeX
  \usepackage[T1]{fontenc}
  \usepackage[utf8]{inputenc}
  \usepackage{textcomp} % provide euro and other symbols
\else % if luatex or xetex
  \usepackage{unicode-math} % this also loads fontspec
  \defaultfontfeatures{Scale=MatchLowercase}
  \defaultfontfeatures[\rmfamily]{Ligatures=TeX,Scale=1}
\fi
\usepackage{lmodern}
\ifPDFTeX\else
  % xetex/luatex font selection
\fi
% Use upquote if available, for straight quotes in verbatim environments
\IfFileExists{upquote.sty}{\usepackage{upquote}}{}
\IfFileExists{microtype.sty}{% use microtype if available
  \usepackage[]{microtype}
  \UseMicrotypeSet[protrusion]{basicmath} % disable protrusion for tt fonts
}{}
\makeatletter
\@ifundefined{KOMAClassName}{% if non-KOMA class
  \IfFileExists{parskip.sty}{%
    \usepackage{parskip}
  }{% else
    \setlength{\parindent}{0pt}
    \setlength{\parskip}{6pt plus 2pt minus 1pt}}
}{% if KOMA class
  \KOMAoptions{parskip=half}}
\makeatother
\usepackage{longtable,booktabs,array}
\newcounter{none} % for unnumbered tables
\usepackage{calc} % for calculating minipage widths
\usepackage{caption}
% Make caption package work with longtable
\makeatletter
\def\fnum@table{\tablename~\thetable}
\makeatother
\setlength{\emergencystretch}{3em} % prevent overfull lines
\providecommand{\tightlist}{%
  \setlength{\itemsep}{0pt}\setlength{\parskip}{0pt}}
\usepackage{bookmark}
\IfFileExists{xurl.sty}{\usepackage{xurl}}{} % add URL line breaks if available
\urlstyle{same}
\hypersetup{
  hidelinks,
  pdfcreator={LaTeX via pandoc}}

\author{\texorpdfstring{}{}}
\date{}

\begin{document}

\begin{frame}{The Sound Money Rebellion: Re-Monetizing Tangible
Sovereignty (1985--2026)}
\protect\phantomsection\label{the-sound-money-rebellion-re-monetizing-tangible-sovereignty-19852026}
Yet the resistance to fiat abstraction has proven structurally
persistent. Despite the closure of the gold window in 1971, the demand
for tangible, unmediated value has periodically erupted with
force---most dramatically in the Hunt brothers' attempted silver
accumulation (1979--80) and the subsequent establishment of the American
Gold Eagle (1985) and American Silver Eagle (1986) bullion programs,
which reintroduced sovereign metallic bullion not for circulation but as
constitutionally significant repositories of stored value. By January
2026, the gold-to-silver ratio had compressed from the historical range
of 80--90:1 that prevailed through most of the 2010s toward
approximately 46:1, suggesting a structural reassertion of silver's
monetary premium alongside its growing industrial demand profile.

\begin{block}{The Legislative Renaissance of American State Monetary
Federalism}
\protect\phantomsection\label{the-legislative-renaissance-of-american-state-monetary-federalism}
The battle between ``allegoric riches'' and metallic tangibility has
escalated directly into the 21st century as an assertion of
constitutional sovereignty. Beginning with the \textbf{Utah Legal Tender
Act} (HB 317, signed 25 March 2011), a coalition of US states has
embarked on a coordinated legislative campaign to formally challenge
federal monetary hegemony {[}@utah2011legal;
@greene2015understanding{]}. This movement asserts that constituent
states possess the constitutional authority under \textbf{Article I,
Section 10} of the US Constitution---which specifically commands that no
state shall ``make any Thing but gold and silver Coin a Tender in
Payment of Debts''---to erect a parallel monetary ecosystem protecting
intra-state commerce from the depreciating trajectory of the fiat
dollar.

This constitutional clause is not a footnote; it is the enduring
legislative trace of a vision Blake would have recognized. The Framers
of the 1787 Constitution---working two years before Blake etched
\emph{The Book of Thel} with its metallic interrogative---inscribed into
the supreme law of the republic a mandate that state-level money must be
physically grounded. Article I, Section 10 compels a monetary realism
that the subsequent history of fiat abstraction has sought to render
operationally moot without legally repealing. The contemporary sound
money movement is, in this constitutional reading, not an innovation but
a \emph{re-activation} of dormant foundational law---what Blake's Los,
insisting on the tangible and the particular against the geometric
abstraction of the ledger, would recognize as the reassertion of the
Forge against the Mint.

Between 2011 and 2026, this movement has accelerated with legislative
momentum. As of March 2026, \textbf{at least ten states} have passed
legislation formally operationalizing this constitutional clause:

{\def\LTcaptype{none} % do not increment counter
\begin{longtable}[]{@{}
  >{\raggedright\arraybackslash}p{(\linewidth - 6\tabcolsep) * \real{0.1750}}
  >{\raggedright\arraybackslash}p{(\linewidth - 6\tabcolsep) * \real{0.2750}}
  >{\raggedright\arraybackslash}p{(\linewidth - 6\tabcolsep) * \real{0.1500}}
  >{\raggedright\arraybackslash}p{(\linewidth - 6\tabcolsep) * \real{0.4000}}@{}}
\toprule\noalign{}
\begin{minipage}[b]{\linewidth}\raggedright
State
\end{minipage} & \begin{minipage}[b]{\linewidth}\raggedright
Legislation
\end{minipage} & \begin{minipage}[b]{\linewidth}\raggedright
Year
\end{minipage} & \begin{minipage}[b]{\linewidth}\raggedright
Key Provisions
\end{minipage} \\
\midrule\noalign{}
\endhead
State & Legislation & Year & Key Provisions \\
------- & ----------- & ------ & ---------------- \\
Utah & HB 317 & 2011 & First state to recognize US-minted gold/silver
coins as legal tender {[}@utah2011legal{]} \\
Oklahoma & SB 862 & 2014 & Legal tender for US-minted coins; full tax
exemption \\
Arizona & HB 2014 & 2017 & Sales tax and capital gains tax exemption on
precious metals \\
Louisiana & HB 396 / SB 232 & 2017/2024 & Sales tax exemption;
reaffirmed legal tender for US coins \\
Wyoming & HB 103 & 2018 & Legal tender recognition; sales tax exempt
{[}@wyoming2018legaltender{]} \\
Arkansas & Act 595 & 2023 & Full legal tender recognition and tax
exemption {[}@arkansas2023legaltender{]} \\
Florida & CS/HB 999 & 2025 & Legal tender for payment of debts;
effective July 1, 2026 {[}@florida2025legaltender{]} \\
Texas & HB 1056 & 2025 & Signed by Governor Abbott June 2025; partial
effect Sept 2026 {[}@texas2025legaltender{]} \\
Alabama & SB 130 & 2025 & Passed House 102-0; comprehensive legal tender
and tax provisions {[}@alabama2025legaltender{]} \\
Nebraska & LB 1317 & 2024 & Removes capital gains tax; explicitly
excludes CBDCs from definition of money {[}@nebraska2024metals{]} \\
\bottomrule\noalign{}
\end{longtable}
}

This is not fringe activism. The legislative margins---Alabama's SB 130
passing the House 102-0---indicate broad bipartisan consensus that the
constitutional prerogative of metallic money has been temporarily
suspended rather than lawfully repealed. The Sound Money Defense League
reports broad state-level movement toward reduced taxation on precious
metals, representing a national trend toward the formal re-recognition
of gold and silver's constitutional monetary status.
\end{block}

\begin{block}{Dismantling the Punitive Tax Architecture of Fiat
Enforcement}
\protect\phantomsection\label{dismantling-the-punitive-tax-architecture-of-fiat-enforcement}
A critical dimension of this monetary federalism involves the systematic
dismantling of the tax architecture designed to penalize precious metal
ownership. Under federal tax law, the IRS classifies bullion as a
``collectible'' subject to a 28\% capital gains rate---significantly
higher than the 15--20\% rate applied to stocks, bonds, and real
estate---structurally discouraging its use as a savings vehicle or
transactional medium. This classification uses the tax code to enforce
fiat monopoly: any citizen who seeks to preserve purchasing power in
constitutional money is penalized relative to those who hold Federal
Reserve Notes.

The state-level response has been to systematically strip away these
penalties. Utah's pioneering HB 317 (2011), later joined by additional
measures expanding legal-tender recognition and payment infrastructure,
has been followed by Arizona and Oklahoma's complete exemption of sales
taxes and capital gains on gold and silver; Nebraska's elimination of
state capital gains taxes on metals; and Florida and Alabama's abolition
of sales taxes on bullion acquisitions {[}@utah2011legal;
@florida2025legaltender; @alabama2025legaltender;
@nebraska2024metals{]}. Nebraska's LB 1317 is particularly significant:
in addition to removing capital gains taxes on precious metals, it
explicitly \textbf{excludes Central Bank Digital Currencies (CBDCs) from
the legal definition of money}---a direct legislative assertion that
digital fiat tokens issued by the Federal Reserve do not constitute
constitutionally valid money {[}@nebraska2024metals{]}.

By eliminating the tax friction designed to penalize the holding of
physical metal, these states directly challenge the Federal Reserve's
capacity to define and administer the medium of exchange. In the
language of \emph{America a Prophecy}, this is the Thirteen Angels'
collective transformation writ legislative: the states, like Blake's
revolutionary governors who rend their robes and stand with Washington
in the Orcian flames, are breaking with the received institutional
framework---enacting the pragmatist moment of genuine inquiry, where
inherited habit is abandoned in favor of a new mode of engagement with
the indeterminate monetary situation {[}@friedman2026blake{]}. The
movement explicitly positions itself as a defense of the producing
classes---an insistence that constitutional metals, constrained by
physical scarcity, preserve the purchasing power of labor against the
indefinite expansion of fiat currency. It is the modern institutional
manifestation of Bryan's agrarian crusade, transposed from the prairie
to the statehouse and equipped with precise statutory instruments.
\end{block}

\begin{block}{The Texas Bullion Depository: Sovereign Infrastructure
Beyond Federal Custody}
\protect\phantomsection\label{the-texas-bullion-depository-sovereign-infrastructure-beyond-federal-custody}
A key structural innovation of this 21st-century resistance is the
establishment of state bullion depositories. In 2015, the Texas
legislature passed \textbf{SB 2097}, creating the \textbf{Texas Bullion
Depository}---launched operationally in 2018 as the first sovereign
institutional infrastructure designed specifically for holding precious
metals reserves outside of the federal banking apparatus
{[}@texas2015depository{]}. The Depository provides a secure,
state-managed, publicly auditable alternative to the federal
depositories (Fort Knox has not been independently audited since 1953),
and is authorized to accept deposits from individuals, businesses, and
government entities. Governor Abbott's subsequent signing of \textbf{HB
1056} (June 2025) integrates the Depository's infrastructure with the
state's new legal tender framework, creating a vertically integrated
system: Texas citizens can hold gold and silver as legal tender, store
it in a state-managed depository, and transact with it free of state
taxation {[}@texas2025legaltender{]}.

This infrastructure physically relocates monetary reserves from
centralized Federal Reserve custody into explicitly decentralized,
state-controlled institutions. When mapped onto Blake's prophetic
architecture, this localized legislative reassertion represents the
persistence of what Blake called the productive, creative impulse of
Los---the blacksmith who works with fire and metal at the forge,
insisting on the tangible and the particular against the geometric
abstraction of the ledger. The constitutional text that these states
invoke---Article I, Section 10---was drafted in 1787, two years before
\emph{The Book of Thel} (1789) {[}@blake1789thel{]}. Over two centuries
later, the tension between abstract ledger and tangible metal endures,
and the legislative campaign to reclaim physical monetary sovereignty
remains a persistent, structural defense of tangible value against the
empire of abstract debt {[}@erdman1988complete{]}.
\end{block}
\end{frame}

\end{document}
